\section{Evaluation Design}\label{sec:eval}

\ins{Explain how you evaluated your approach and why your evaluation is appropriate for the problem. Describe datasets, metrics, baselines, and comparisons. The emphasis is on sound experimental reasoning, not on outperforming prior work.}

% % Experiments, Results & Analysis
% % 5
% % Appropriate evaluation design, meaningful results, thoughtful analysis and interpretation, and discussion of limitations or failure cases, and implications for future work.

% Our evaluation framework measures two complementary dimensions: \emph{susc
% ceptibility} and \emph{transparency}.

% \paragraph{Behavioral susceptibility.}
% We first determine whether the injected hint changes the model’s answer relative to the no-hint baseline. If the model changes from a correct baseline answer to an incorrect hinted answer, we classify the behavior as \emph{sycophantic} or socially influenced.

% \paragraph{Reasoning transparency.}
% We then evaluate whether the model explicitly acknowledges the hint within its CoT explanation. We detect hint acknowledgment using a combination of keyword matching and LLM-as-a-judge evaluation. For example, authority hints are identified through references to ``professor'' or ``expert,'' while social hints are identified through references to group consensus.

% \paragraph{Faithfulness labels.}
% Using these two axes, we categorize model behavior into three outcomes:

% \begin{itemize}
%     \item \textit{Resisted:} The answer does not change under hint injection.
    
%     \item \textit{Transparent Sycophancy:} The answer changes and the reasoning trace explicitly references the hint.
    
%     \item \textit{Opaque Sycophancy:} The answer changes but the reasoning trace does not acknowledge the hint.
% \end{itemize}

% We treat opaque sycophancy as the strongest failure of reasoning faithfulness because the model’s behavior is influenced by the injected cue without transparently reporting that influence in its explanation.


Our evaluation framework is designed to measure whether reasoning models are influenced by socially grounded bias cues and whether those influences are transparently reflected in generated chain-of-thought (CoT) explanations. Because our goal is to evaluate reasoning faithfulness rather than benchmark accuracy, we focus on comparing model behavior across hint types, language conditions, and task domains.

\subsection{Faithfulness Metrics}

Our evaluation framework measures two complementary dimensions: \emph{behavioral susceptibility} and \emph{reasoning transparency}.

\paragraph{Behavioral susceptibility.}
We first determine whether the injected hint changes the model’s answer relative to the no-hint baseline. If the model changes from a correct baseline answer to an incorrect hinted answer, we classify the behavior as \emph{sycophantic} or socially influenced.

\paragraph{Reasoning transparency.}
We then evaluate whether the model explicitly acknowledges the hint within its CoT explanation. We detect hint acknowledgment using a combination of keyword matching and LLM-as-a-judge evaluation. For example, authority hints are identified through references to ``professor'' or ``expert,'' while social hints are identified through references to group consensus.

\paragraph{Faithfulness labels.}
Using these two axes, we categorize model behavior into three outcomes:

\begin{itemize}
    \item \textit{Resisted:} The answer does not change under hint injection.
    
    \item \textit{Transparent Sycophancy:} The answer changes and the reasoning trace explicitly references the hint.
    
    \item \textit{Opaque Sycophancy:} The answer changes but the reasoning trace does not acknowledge the hint.
\end{itemize}

We treat opaque sycophancy as the strongest failure of reasoning faithfulness because the model’s behavior is influenced by the injected cue without transparently reporting that influence in its explanation.

\subsection{Comparative Analysis}
We compare different categories of hint injection (authority, social, and indirect) in order to evaluate whether some forms of social influence are more likely to produce unfaithful or opaque reasoning behavior. We hypothesize that indirect hints may be less likely to appear in reasoning traces because they provide weaker or less explicit social signals.

To evaluate whether cross-lingual mismatch affects either behavioral susceptibility or reasoning transparency, we compare matched and mismatched language conditions. This analysis directly tests whether multilingual context changes the likelihood that models explicitly acknowledge socially grounded influence in their explanations.

Then, we compare culturally grounded reasoning tasks against mathematical reasoning tasks. Cultural tasks derived from CultureBank involve implicit social norms and contextual interpretation, while GSM8K math tasks provide a more objective reasoning setting with verifiable answers. This comparison allows us to determine whether observed faithfulness failures are specific to culturally grounded reasoning or generalize across domains.

Finally, because some reasoning-enabled outputs do not contain explicit CoT traces, we additionally analyze the frequency of missing or suppressed reasoning across conditions. We treat these cases as potentially meaningful model behavior rather than simple generation failures, since the absence of reasoning may itself reflect differences in how models process ambiguous or socially grounded hints.
